% 矩阵分解示意；原图使用W标签，正文用G，保留差异。
\begin{figure}[htbp]
\centering
\begin{tikzpicture}[x=.62cm,y=.62cm,font=\small,>=Stealth]
\definecolor{chSixGreen}{RGB}{0,167,80}
\definecolor{chSixBlue}{RGB}{52,104,159}
\fill[blue!18] (0,0) rectangle (6,6);
\draw[chSixBlue] (0,0) grid[step=1] (6,6);
\draw[chSixBlue,thick] (0,0) rectangle (6,6);
\fill[blue!18] (7.6,0) rectangle (11.6,6);
\draw[chSixBlue,xstep=1,ystep=1] (7.6,0) grid (11.6,6);
\draw[chSixBlue,thick] (7.6,0) rectangle (11.6,6);
\fill[blue!18] (13.1,3.6) rectangle (19.1,6);
\foreach \xx in {13.1,14.1,...,19.1} \draw[chSixBlue] (\xx,3.6)--(\xx,6);
\foreach \yy in {3.6,4.4,5.2,6} \draw[chSixBlue] (13.1,\yy)--(19.1,\yy);
\draw[chSixBlue,thick] (13.1,3.6) rectangle (19.1,6);
\draw[->,chSixGreen,line width=4pt] (6.25,3)--(7.3,3);
\node[red] at (3,3) {$X$};
\node[red] at (9.6,3) {$W$};
\node[red] at (16.1,4.8) {$H$};
\node[anchor=north east] at (0,0) {$m$};
\node[anchor=north east] at (7.6,0) {$m$};
\node[anchor=south east] at (6,6) {$n$};
\node[anchor=south west] at (11.6,6) {$r$};
\node[anchor=south east] at (19.1,6) {$n$};
\node[anchor=north east] at (13.1,3.6) {$r$};
\draw[->,chSixGreen] (-1.4,3)--(0,3);
\node[chSixGreen,anchor=east] at (-.1,2.6) {特征};
\draw[->,chSixGreen] (.7,7.1)--(.7,6);
\node[chSixGreen,anchor=west] at (.8,6.6) {体素（数据点）};
\draw[->,chSixGreen] (8,7.8)--(8,6);
\node[chSixGreen,align=left,anchor=west] at (8.1,7.2) {子区域／\\组织类型};
\draw[->,red] (8.9,6.5)--(12.5,6.5)--(12.5,5.2)--(13.1,5.2);
\node[red,anchor=west] at (12.5,6.2) {比例};
\end{tikzpicture}
\origfig{1}{矩阵分解。}
\end{figure}
